% !TEX root = marti_ferran_gonzales_cv.tex
% Engine is chosen in .latexmkrc. Do not add a TeX-program magic comment:
% the editor would run the engine directly and dump auxiliary files here.
\documentclass[a4paper,11pt]{article}
\newcommand{\cvsetupfont}{\setsansfont{Calibri}}
% Shared layout for marti_ferran_gonzales_cv.tex and marti_ferran_gonzales_cover_letter.tex.
% The CV defines \cvsetupfont before % Shared layout for marti_ferran_gonzales_cv.tex and marti_ferran_gonzales_cover_letter.tex.
% The CV defines \cvsetupfont before % Shared layout for marti_ferran_gonzales_cv.tex and marti_ferran_gonzales_cover_letter.tex.
% The CV defines \cvsetupfont before \input{preamble}.
% Author : Jake Gutierrez
% Based off of: https://github.com/sb2nov/resume
% License : MIT
% This file is not a document. Compile marti_ferran_gonzales_cv.tex or marti_ferran_gonzales_cover_letter.tex.

\ifdefined\cvsetupfont
  \usepackage{fontspec}
  \cvsetupfont
  \renewcommand{\familydefault}{\sfdefault}
\else
  \usepackage[T1]{fontenc}
  \usepackage[utf8]{inputenc}
\fi
\usepackage{latexsym}
\usepackage[empty]{fullpage}
\usepackage{titlesec}
\usepackage{marvosym}
\usepackage[usenames,dvipsnames]{color}
\usepackage{verbatim}
\usepackage{enumitem}
\usepackage[hidelinks]{hyperref}
\usepackage{fancyhdr}
\usepackage[english]{babel}
\usepackage{tabularx}
\ifdefined\cvsetupfont
\else
  \input{glyphtounicode}
\fi

\pagestyle{fancy}
\fancyhf{} % clear all header and footer fields
\fancyfoot{}
\renewcommand{\headrulewidth}{0pt}
\renewcommand{\footrulewidth}{0pt}

% Adjust margins
\addtolength{\oddsidemargin}{-0.5in}
\addtolength{\evensidemargin}{-0.5in}
\addtolength{\textwidth}{1in}
\addtolength{\topmargin}{-.5in}
\addtolength{\textheight}{1.0in}

\urlstyle{same}

\raggedbottom
\raggedright
\setlength{\tabcolsep}{0in}

% Sans faces here have no small caps, so the section title stays in the same font.
\ifdefined\cvsetupfont
  \titleformat{\section}{
    \vspace{-4pt}\raggedright\large
  }{}{0em}{\MakeUppercase}[\color{black}\titlerule \vspace{-5pt}]
\else
  \titleformat{\section}{
    \vspace{-4pt}\scshape\raggedright\large
  }{}{0em}{}[\color{black}\titlerule \vspace{-5pt}]
  % Ensure that the generated pdf is machine readable/ATS parsable
  \pdfgentounicode=1
\fi

%-------------------------
% Custom commands
\newcommand{\resumeItem}[1]{
  \item\small{
    {#1 \vspace{-2pt}}
  }
}

\newcommand{\resumeSubheading}[4]{
  \vspace{-2pt}\item
    \begin{tabular*}{0.97\textwidth}[t]{l@{\extracolsep{\fill}}r}
      \textbf{#1} & #2 \\
      \textit{\small#3} & \textit{\small #4} \\
    \end{tabular*}\vspace{-7pt}
}

\newcommand{\resumeSubSubheading}[2]{
    \item
    \begin{tabular*}{0.97\textwidth}{l@{\extracolsep{\fill}}r}
      \textit{\small#1} & \textit{\small #2} \\
    \end{tabular*}\vspace{-7pt}
}

\newcommand{\resumeProjectHeading}[2]{
    \item
    \begin{tabular*}{0.97\textwidth}{l@{\extracolsep{\fill}}r}
      \small#1 & #2 \\
    \end{tabular*}\vspace{-7pt}
}

\newcommand{\resumeSubItem}[1]{\resumeItem{#1}\vspace{-4pt}}

\renewcommand\labelitemii{$\vcenter{\hbox{\tiny$\bullet$}}$}

\newcommand{\resumeSubHeadingListStart}{\begin{itemize}[leftmargin=0.15in, label={}]}
\newcommand{\resumeSubHeadingListEnd}{\end{itemize}}
\newcommand{\resumeItemListStart}{\begin{itemize}}
\newcommand{\resumeItemListEnd}{\end{itemize}\vspace{-5pt}}

% Role context under a \resumeSubheading. The heading table is top-aligned and
% pulls the next line up, so the blurb must start a new paragraph below it.
\newcommand{\resumeIntro}[1]{\par\vspace{3pt}{\small #1\par}}
.
% Author : Jake Gutierrez
% Based off of: https://github.com/sb2nov/resume
% License : MIT
% This file is not a document. Compile marti_ferran_gonzales_cv.tex or marti_ferran_gonzales_cover_letter.tex.

\ifdefined\cvsetupfont
  \usepackage{fontspec}
  \cvsetupfont
  \renewcommand{\familydefault}{\sfdefault}
\else
  \usepackage[T1]{fontenc}
  \usepackage[utf8]{inputenc}
\fi
\usepackage{latexsym}
\usepackage[empty]{fullpage}
\usepackage{titlesec}
\usepackage{marvosym}
\usepackage[usenames,dvipsnames]{color}
\usepackage{verbatim}
\usepackage{enumitem}
\usepackage[hidelinks]{hyperref}
\usepackage{fancyhdr}
\usepackage[english]{babel}
\usepackage{tabularx}
\ifdefined\cvsetupfont
\else
  \input{glyphtounicode}
\fi

\pagestyle{fancy}
\fancyhf{} % clear all header and footer fields
\fancyfoot{}
\renewcommand{\headrulewidth}{0pt}
\renewcommand{\footrulewidth}{0pt}

% Adjust margins
\addtolength{\oddsidemargin}{-0.5in}
\addtolength{\evensidemargin}{-0.5in}
\addtolength{\textwidth}{1in}
\addtolength{\topmargin}{-.5in}
\addtolength{\textheight}{1.0in}

\urlstyle{same}

\raggedbottom
\raggedright
\setlength{\tabcolsep}{0in}

% Sans faces here have no small caps, so the section title stays in the same font.
\ifdefined\cvsetupfont
  \titleformat{\section}{
    \vspace{-4pt}\raggedright\large
  }{}{0em}{\MakeUppercase}[\color{black}\titlerule \vspace{-5pt}]
\else
  \titleformat{\section}{
    \vspace{-4pt}\scshape\raggedright\large
  }{}{0em}{}[\color{black}\titlerule \vspace{-5pt}]
  % Ensure that the generated pdf is machine readable/ATS parsable
  \pdfgentounicode=1
\fi

%-------------------------
% Custom commands
\newcommand{\resumeItem}[1]{
  \item\small{
    {#1 \vspace{-2pt}}
  }
}

\newcommand{\resumeSubheading}[4]{
  \vspace{-2pt}\item
    \begin{tabular*}{0.97\textwidth}[t]{l@{\extracolsep{\fill}}r}
      \textbf{#1} & #2 \\
      \textit{\small#3} & \textit{\small #4} \\
    \end{tabular*}\vspace{-7pt}
}

\newcommand{\resumeSubSubheading}[2]{
    \item
    \begin{tabular*}{0.97\textwidth}{l@{\extracolsep{\fill}}r}
      \textit{\small#1} & \textit{\small #2} \\
    \end{tabular*}\vspace{-7pt}
}

\newcommand{\resumeProjectHeading}[2]{
    \item
    \begin{tabular*}{0.97\textwidth}{l@{\extracolsep{\fill}}r}
      \small#1 & #2 \\
    \end{tabular*}\vspace{-7pt}
}

\newcommand{\resumeSubItem}[1]{\resumeItem{#1}\vspace{-4pt}}

\renewcommand\labelitemii{$\vcenter{\hbox{\tiny$\bullet$}}$}

\newcommand{\resumeSubHeadingListStart}{\begin{itemize}[leftmargin=0.15in, label={}]}
\newcommand{\resumeSubHeadingListEnd}{\end{itemize}}
\newcommand{\resumeItemListStart}{\begin{itemize}}
\newcommand{\resumeItemListEnd}{\end{itemize}\vspace{-5pt}}

% Role context under a \resumeSubheading. The heading table is top-aligned and
% pulls the next line up, so the blurb must start a new paragraph below it.
\newcommand{\resumeIntro}[1]{\par\vspace{3pt}{\small #1\par}}
.
% Author : Jake Gutierrez
% Based off of: https://github.com/sb2nov/resume
% License : MIT
% This file is not a document. Compile marti_ferran_gonzales_cv.tex or marti_ferran_gonzales_cover_letter.tex.

\ifdefined\cvsetupfont
  \usepackage{fontspec}
  \cvsetupfont
  \renewcommand{\familydefault}{\sfdefault}
\else
  \usepackage[T1]{fontenc}
  \usepackage[utf8]{inputenc}
\fi
\usepackage{latexsym}
\usepackage[empty]{fullpage}
\usepackage{titlesec}
\usepackage{marvosym}
\usepackage[usenames,dvipsnames]{color}
\usepackage{verbatim}
\usepackage{enumitem}
\usepackage[hidelinks]{hyperref}
\usepackage{fancyhdr}
\usepackage[english]{babel}
\usepackage{tabularx}
\ifdefined\cvsetupfont
\else
  \input{glyphtounicode}
\fi

\pagestyle{fancy}
\fancyhf{} % clear all header and footer fields
\fancyfoot{}
\renewcommand{\headrulewidth}{0pt}
\renewcommand{\footrulewidth}{0pt}

% Adjust margins
\addtolength{\oddsidemargin}{-0.5in}
\addtolength{\evensidemargin}{-0.5in}
\addtolength{\textwidth}{1in}
\addtolength{\topmargin}{-.5in}
\addtolength{\textheight}{1.0in}

\urlstyle{same}

\raggedbottom
\raggedright
\setlength{\tabcolsep}{0in}

% Sans faces here have no small caps, so the section title stays in the same font.
\ifdefined\cvsetupfont
  \titleformat{\section}{
    \vspace{-4pt}\raggedright\large
  }{}{0em}{\MakeUppercase}[\color{black}\titlerule \vspace{-5pt}]
\else
  \titleformat{\section}{
    \vspace{-4pt}\scshape\raggedright\large
  }{}{0em}{}[\color{black}\titlerule \vspace{-5pt}]
  % Ensure that the generated pdf is machine readable/ATS parsable
  \pdfgentounicode=1
\fi

%-------------------------
% Custom commands
\newcommand{\resumeItem}[1]{
  \item\small{
    {#1 \vspace{-2pt}}
  }
}

\newcommand{\resumeSubheading}[4]{
  \vspace{-2pt}\item
    \begin{tabular*}{0.97\textwidth}[t]{l@{\extracolsep{\fill}}r}
      \textbf{#1} & #2 \\
      \textit{\small#3} & \textit{\small #4} \\
    \end{tabular*}\vspace{-7pt}
}

\newcommand{\resumeSubSubheading}[2]{
    \item
    \begin{tabular*}{0.97\textwidth}{l@{\extracolsep{\fill}}r}
      \textit{\small#1} & \textit{\small #2} \\
    \end{tabular*}\vspace{-7pt}
}

\newcommand{\resumeProjectHeading}[2]{
    \item
    \begin{tabular*}{0.97\textwidth}{l@{\extracolsep{\fill}}r}
      \small#1 & #2 \\
    \end{tabular*}\vspace{-7pt}
}

\newcommand{\resumeSubItem}[1]{\resumeItem{#1}\vspace{-4pt}}

\renewcommand\labelitemii{$\vcenter{\hbox{\tiny$\bullet$}}$}

\newcommand{\resumeSubHeadingListStart}{\begin{itemize}[leftmargin=0.15in, label={}]}
\newcommand{\resumeSubHeadingListEnd}{\end{itemize}}
\newcommand{\resumeItemListStart}{\begin{itemize}}
\newcommand{\resumeItemListEnd}{\end{itemize}\vspace{-5pt}}

% Role context under a \resumeSubheading. The heading table is top-aligned and
% pulls the next line up, so the blurb must start a new paragraph below it.
\newcommand{\resumeIntro}[1]{\par\vspace{3pt}{\small #1\par}}


\begin{document}
\begin{center}
    \textbf{\Huge MARTÍ FERRAN GONZALES} \\ \vspace{1pt}
    \small Cranfield, United Kingdom -- Barcelona, Spain \\ \vspace{1pt}
    \small +44 5673 224450 -- marti.ferran@icloud.com -- linkedin.com/in/marti-ferran-gonzales
\end{center}

\section{Personal Statement}
{\small Practical and analytical aerospace engineer currently completing a double-degree master’s programme at Cranfield University and UPC, specialising in autonomous vehicle dynamics, modelling, and control systems. Brings four years of hands-on experience in Formula Student and professional motorsport at CUPRA Racing, connecting modelling and simulation with on-track prototype testing. A collaborative team player and technical leader who is comfortable in fast-paced environments, coordinating sub-teams and communicating fluently in English, Spanish, and Catalan.}

\section{Key Achievements}
  \resumeItemListStart
    \resumeItem{Achieved a mark of 10/10 for a bachelor's thesis on multi-physics vehicle modelling and the design and testing of a closed-loop yaw stability controller, developed during an internship at CUPRA Racing (CUPRA Racing and UPC, 2025).}
    \resumeItem{Refactored the VCU software architecture in Simulink, reorganising timer management and task scheduling to eliminate missed driver inputs and false CAN timeouts, improve execution speed and ensure consistent CAN message timing (CUPRA Racing, 2026).}
    \resumeItem{Helped UPC ecoRacing achieve 2nd place overall at Formula Student Portugal and 3rd place overall at Formula Student Spain, contributing to vehicle dynamics, control software development and track testing (UPC ecoRacing, 2023).}
    \resumeItem{Awarded an honours distinction for achieving the highest grade in Airport Design and Construction during the MSc in Aeronautical Engineering (Universitat Politècnica de Catalunya, 2026).}
  \resumeItemListEnd

\section{Education}
  \resumeSubHeadingListStart
    \resumeSubheading
      {MSc in Autonomous Vehicle Dynamics and Control}{September 2026 -- August 2027}
      {Cranfield University (Cranfield, UK)}{}
      \resumeItemListStart
        \resumeItem{Modules: Guidance and Navigation for Autonomous Systems; Sensor Fusion; UAS Dynamics and Control; UAS Modelling and Simulation; Autonomous Vehicle Control Systems; AI for Autonomous Systems; Aerial Communication Systems}
      \resumeItemListEnd
    \resumeSubheading
      {MSc in Aeronautical Engineering}{September 2025 -- August 2027}
      {Universitat Politècnica de Catalunya (Terrassa, Spain)}{}
      \resumeItemListStart
        \resumeItem{Current Average Grade: 8.83/10 (60 ECTS completed; 120 ECTS total; Cranfield double degree)}
        \resumeItem{Modules: Aerodynamics, Flight and Orbital Mechanics; Air Transport and Navigation Systems; Aerospace Vehicles; RF and Communication Systems; Aircraft and Rocket Propulsion Systems; Airport Design and Construction (Honours Distinction, 2026)}
      \resumeItemListEnd
    \resumeSubheading
      {BSc in Aerospace Technologies Engineering}{September 2021 -- September 2025}
      {Universitat Politècnica de Catalunya (Terrassa, Spain)}{}
      \resumeItemListStart
        \resumeItem{Average Grade: 7.64/10 (240 ECTS)}
        \resumeItem{Modules: Automatic Control; Flight Mechanics; Avionics; Electronic Circuits; Computational Engineering; Aerodynamics; Aerospace Structures; Aerospace Propulsion; Space Engineering; Aircraft Design}
        \resumeItem{Bachelor's Thesis: Study of yaw moment and stability control algorithm based on longitudinal torque distribution in electric vehicles --- Grade: 10/10 (Nominated for Honours Distinction)}
      \resumeItemListEnd
  \resumeSubHeadingListEnd

\section{Career History}
  \resumeSubHeadingListStart
    \resumeSubheading
      {Electronics and Embedded Software Engineering Intern}{January 2025 -- February 2026}
      {CUPRA Racing (Abrera, Spain)}{}
      \resumeIntro{Motorsport division of CUPRA, responsible for the development of high-performance racing vehicles, experimental prototype platforms, and advanced engineering R\&D projects.}
      \resumeItemListStart
        \resumeItem{Developed a complete multi-physics vehicle simulation model from scratch in MATLAB/Simulink for the TCR Hybrid prototype, incorporating ICE test-bench data, tyre and suspension dynamics, and used it to test and calibrate yaw control algorithms based on longitudinal torque distribution before track testing.}
        \resumeItem{Developed closed-loop yaw control algorithms for three prototypes using torque distribution, including independent left-right torque control through in-wheel motors on one prototype. Calibrated and validated the algorithms through track testing with professional drivers, with telemetry and driver feedback confirming improvements in vehicle balance, stability and control.}
        \resumeItem{Completely refactored the TCR Hybrid VCU software in Simulink, reorganising timer management and task scheduling. This eliminated missed driver inputs and false CAN timeouts, improved execution speed, and made the timing of periodic tasks, such as light flashing, more consistent.}
        \resumeItem{Developed thermal models of the TCR Hybrid battery and electric motor in MATLAB using existing empirical system data, then simulated race conditions to determine coolant flow rates, size radiators and establish sustainable battery and motor power limits for a full race.}
        \resumeItem{Configured dataloggers and acquisition tables for prototype track testing, including sensor selection, and analysed telemetry to diagnose faults and perform control system validation and calibration.}
      \resumeItemListEnd

    \resumeSubheading
      {Vehicle Dynamics \& Control Engineer}{November 2021 -- September 2024}
      {UPC ecoRacing (Formula Student) (Terrassa, Spain)}{}
      \resumeIntro{Spanish Formula Student team, designing and manufacturing high-performance 4WD electric and autonomous (manual and driverless) formula prototypes.}
      \resumeItemListStart
        \resumeItem{Developed a closed-loop torque vectoring algorithm for a 4WD electric prototype with four independent in-wheel motors, improving yaw control robustness and vehicle performance through on-track calibration.}
        \resumeItem{Led the concept design of the 2025 Formula Student prototype, coordinating engineering teams to define the vehicle architecture, make key concept design decisions and set performance targets.}
        \resumeItem{Fitted Magic Formula tyre models to experimental tyre test-rig data for tyre selection, then integrated them into non-linear vehicle models in MATLAB and Python to predict vehicle behaviour and evaluate control algorithms and mechanical setup changes.}
        \resumeItem{Led track testing sessions, planning tests and coordinating drivers and engineers. Used telemetry to validate simulation models and control algorithms, resolve electrical and mechanical faults, and improve on-track performance through setup changes and controller calibration.}
        \resumeItem{Designed and sized suspension and steering components using CAD and FEA to minimise weight, then checked the calculated loads against strain-gauge measurements from track testing.}
      \resumeItemListEnd

    \resumeSubheading
      {Head of Logistics}{September 2023 -- September 2024}
      {UPC ecoRacing (Formula Student) (Terrassa, Spain)}{}
      \resumeItemListStart
        \resumeItem{Managed travel, equipment and inventory for a Formula Student team of over 40 members, coordinating preparations for international competitions.}
        \resumeItem{Worked with engineering sub-teams to organise trackside operations and keep event preparations on schedule.}
      \resumeItemListEnd
  \resumeSubHeadingListEnd

\section{Technical Projects}
  \resumeSubHeadingListStart
    \resumeSubheading
      {Vehicle Dynamics \& Telemetry Analysis Suite}{July 2026 -- Present}
      {System Architecture \& AI-Assisted Development}{}
      \resumeItemListStart
        \resumeItem{Defined the system architecture and engineering logic for a modular desktop software suite focused on telemetry data analysis and suspension kinematics evaluation.}
        \resumeItem{Coordinated development through structured AI-agent workflows to implement data import pipelines, session comparison tools, custom mathematical channels, and visualisation layouts.}
      \resumeItemListEnd
    \resumeSubheading
      {Nutria (Desktop Clinical Management Platform)}{September 2024 -- August 2026}
      {Python, Product Design \& AI Workflows}{}
      \resumeItemListStart
        \resumeItem{Programmed the initial functional prototypes in Python and subsequently led the product architecture and system redesign from concept to a private pilot release.}
        \resumeItem{Defined system requirements and modular architecture, and coordinated AI-assisted development and automated end-to-end testing.}
      \resumeItemListEnd
  \resumeSubHeadingListEnd

\section{Skills, Interests \& Extracurricular Activities}
 \begin{itemize}[leftmargin=0.15in, label={}]
    \small{\item{
     \textbf{Technical Skills:} MATLAB, Simulink, C, C++, Python, LaTeX, SolidWorks, CATIA, FreeCAD, MSC Nastran, ANSA, MoTeC, WinTAX4, Sysma, Git, GitHub, VS Code, Cursor, Antigravity, Claude Code. \\
     \textbf{Languages:} Catalan (Native), Spanish (Native), English (C2 Proficient), German (A1 Basic). \\
     \textbf{Professional/Technical Training:}
     \begin{itemize}
        \item Applied Vehicle Dynamics Seminar --- Optimum G (Stuttgart, November 2024, 48 hours). Certificate of Attendance. Non-linear tyre modelling, vehicle dynamics, kinematics and simulation.
        \item Data-Driven Performance Engineering Seminar --- Optimum G (Stuttgart, November 2024, 48 hours). Certificate of Attendance. Data acquisition, sensor analysis, track testing and performance engineering.
     \end{itemize}
     \textbf{Interests:} Enjoys experimenting with AI tools and exploring new technologies. Outside technology, enjoys spending time with his dogs, staying active through cycling and padel, and looking after plants.
    }}
 \end{itemize}

\end{document}
